% =====================================================================
%  5.8 OBJETIVOS
%  Rubrica: "correctamente definidos y suficientes, y especifican con
%  detalle el sistema". Hay que indicar el ALCANCE de cada uno.
% =====================================================================

\chapter{Objetivos}
\label{cap:objetivos}

\section{Objetivo principal}

Desarrollar una aplicación web de apoyo a la decisión clínica que integre un
motor de inferencia basado en los criterios STOPP/START v3, de modo que el
profesional sanitario pueda revisar de forma automatizada la medicación de
pacientes mayores a partir de los fármacos y diagnósticos del caso.

\section{Objetivos específicos}

\noindent\textbf{OE1: Formalizar los criterios STOPP/START v3 como reglas
evaluables de forma automática.}

\textbf{Alcance:} modelado declarativo de los criterios (STOPP y START) en un
catálogo estructurado, sin sustituir el criterio clínico del profesional.

\medskip
\noindent\textbf{OE2: Implementar un motor de reglas que evalúe el caso
del paciente frente a dichos criterios.}

\textbf{Alcance:} a partir de medicación y diagnósticos, determinar qué avisos
STOPP y START se activan; no incluye integración con historias clínicas
electrónicas externas.

\medskip
\noindent\textbf{OE3: Diseñar un flujo de captura guiado de datos clínicos
relevantes.}

\textbf{Alcance:} interfaz para seleccionar medicamentos y diagnósticos (y,
cuando aplique, datos auxiliares del caso) de forma usable en consulta.

\medskip
\noindent\textbf{OE4: Incorporar ayudas de interfaz que reduzcan el tiempo de
uso.}

\textbf{Alcance:} deshabilitar opciones irrelevantes según la selección actual
y destacar relaciones entre fármacos y diagnósticos; no incluye
recomendaciones terapéuticas libres fuera de STOPP/START.

\medskip
\noindent\textbf{OE5: Presentar los resultados de forma clara y accionable.}

\textbf{Alcance:} visualización de criterios activados (STOPP/START), agrupados
de manera comprensible para la revisión clínica.

\medskip
\noindent\textbf{OE6: Permitir la generación de un informe y la reutilización
del caso.}

\textbf{Alcance:} exportación de informe (p.~ej.\ PDF) y guardado/carga del
caso (p.~ej.\ JSON / persistencia local), sin almacenamiento en servidor.

\medskip
\noindent\textbf{OE7: Validar el sistema mediante pruebas automáticas y
revisión clínica.}

\textbf{Alcance:} batería de tests del motor y de la aplicación, más contraste
con la guía oficial y supervisión de la cotutora clínica; no constituye un
ensayo clínico.

\section{Fuera del alcance del trabajo}

\label{sec:limitaciones}
El presente trabajo se centra en la automatización de la evaluación de los
criterios STOPP/START v3 sobre un caso introducido por el profesional.
Quedan deliberadamente fuera del alcance, entre otros, los siguientes
aspectos:
\begin{itemize}
  \item \textbf{Integración con la historia clínica electrónica.}
  La aplicación no se conecta a sistemas hospitalarios ni de atención
  primaria. En una evolución natural del sistema, la carga del historial
  permitiría preseleccionar automáticamente diagnósticos y medicaciones,
  reduciendo aún más el tiempo de uso en consulta.
  \item \textbf{Sustitución del juicio clínico.}
  La herramienta actúa como apoyo a la decisión: muestra avisos STOPP y
  START aplicables al caso, pero la interpretación y la decisión
  terapéutica final corresponden siempre al profesional sanitario.
  \item \textbf{Catálogo clínico acotado.}
  Medicamentos y diagnósticos disponibles son los modelados en la
  aplicación. No se pretende cubrir la totalidad del arsenal terapéutico
  ni todos los posibles códigos diagnósticos.
  \item \textbf{Arquitectura sin servidor.}
  El sistema opera íntegramente en el navegador, sin backend ni
  almacenamiento centralizado de pacientes, lo que simplifica la
  privacidad en el contexto del TFG pero limita escenarios multiusuario
  o de despliegue institucional.
  \item \textbf{Validación no clínica experimental.}
  La comprobación del sistema se basa en pruebas automáticas y en la
  revisión con la cotutora clínica. No se realiza un ensayo ni un estudio
  de impacto asistencial sobre pacientes reales.
\end{itemize}
